% ============================================================
% main.tex — Documento Raiz do TCC
% Projeto: Alchemon - Jogo Sério para Ensino da Tabela Periódica
% Autor: Mateus Nedel dos Reis
% Formato: IEEE Conference (IEEEtran)
% Compilação: pdflatex -> bibtex -> pdflatex -> pdflatex
%             ou: latexmk -pdf
% ============================================================

\documentclass[conference]{IEEEtran}

% ------------------------------------------------------------
% CODIFICAÇÃO E IDIOMA
% ------------------------------------------------------------
\usepackage[utf8]{inputenc}
\usepackage[T1]{fontenc}
\usepackage[brazil]{babel}

% ------------------------------------------------------------
% GRÁFICOS E MATEMÁTICA
% ------------------------------------------------------------
\usepackage{graphicx}
\graphicspath{{figuras/}}
\usepackage{amsmath}
\usepackage{amssymb}

% ------------------------------------------------------------
% HIPERLINKS E CITAÇÕES
% ------------------------------------------------------------
\usepackage{hyperref}
\hypersetup{
    colorlinks=true,
    linkcolor=black,
    citecolor=black,
    urlcolor=blue,
    pdftitle={Alchemon: Desenvolvimento de um Jogo Sério para o Ensino da Tabela Periódica},
    pdfauthor={Mateus Nedel dos Reis},
    pdfsubject={TCC - Jogos Digitais},
    pdfkeywords={jogos sérios, tabela periódica, gamificação, ensino}
}
\usepackage{cite}

% ------------------------------------------------------------
% TABELAS E LISTAS
% ------------------------------------------------------------
\usepackage{icomma}     % Vírgula decimal em modo matemático (padrão brasileiro)
\usepackage{array}
\usepackage{enumerate}

% ------------------------------------------------------------
% CÓDIGO-FONTE (opcional, para trechos de código)
% ------------------------------------------------------------
\usepackage{listings}
\usepackage{xcolor}

\lstset{
    basicstyle=\ttfamily\footnotesize,
    keywordstyle=\color{blue}\bfseries,
    commentstyle=\color{gray},
    stringstyle=\color{orange},
    breaklines=true,
    frame=single,
    numbers=left,
    numberstyle=\tiny\color{gray},
    captionpos=b
}

% ------------------------------------------------------------
% TÍTULO E AUTORES
% ------------------------------------------------------------
\title{Alchemon: Desenvolvimento de um Jogo Sério\\
para o Ensino da Tabela Periódica}

\author{
    \IEEEauthorblockN{Mateus Nedel dos Reis}
    \IEEEauthorblockA{
        Curso de Tecnologias em Jogos Digitais\\
        Universidade Franciscana\\
        Santa Maria, RS, Brasil\\
        \textit{email@ufn.edu.br}
    }
}

% ============================================================
% INÍCIO DO DOCUMENTO
% ============================================================
\begin{document}

\maketitle
Teste de citacao: \cite{prensky2001digital}
% ------------------------------------------------------------
% RESUMO E PALAVRAS-CHAVE
% ------------------------------------------------------------
\begin{abstract}
Este trabalho apresenta o desenvolvimento do \textit{Alchemon}, um jogo sério
voltado ao ensino da tabela periódica para estudantes do ensino médio. O projeto
explora a interseção entre gamificação e educação em química, propondo uma
solução lúdica e interativa que visa aumentar o engajamento e a retenção do
conteúdo pelos alunos. A abordagem metodológica combina o design de jogos
digitais com princípios pedagógicos modernos, resultando em um produto educativo
avaliado por métricas de aprendizagem e experiência do usuário. Os resultados
preliminares indicam que a mecânica baseada em coleta e batalha de elementos
favorece a memorização das propriedades dos elementos químicos, contribuindo
para um aprendizado mais significativo e duradouro.
\end{abstract}

\begin{IEEEkeywords}
jogos sérios, tabela periódica, gamificação, ensino de química, jogos digitais educacionais
\end{IEEEkeywords}


% ------------------------------------------------------------
% SEÇÕES — cada uma em arquivo separado
% ------------------------------------------------------------
% Chapter 1: Introduction
\section{Introdução}

A introdução serve para apresentar o tema do trabalho, contextualizando o leitor sobre a área de estudo. Aqui, descrevemos a importância do ensino de química, especificamente da Tabela Periódica, e introduzimos o conceito de jogos sérios como uma abordagem pedagógica inovadora.

Este trabalho propõe o desenvolvimento de "Alchemon", um jogo sério projetado para facilitar a aprendizagem dos elementos químicos. A motivação para este projeto vem da observação de que métodos de ensino tradicionais podem não ser suficientemente engajadores para a geração atual de estudantes, que estão imersos em mídias digitais. Como aponta Gee em seu trabalho sobre video games e aprendizagem, ambientes de jogo bem projetados podem ser máquinas de aprendizado poderosas \cite{journal_example}.

O objetivo é criar uma ferramenta que seja ao mesmo tempo educativa e divertida, transformando o estudo da Tabela Periódica em uma experiência interativa e memorável. Acreditamos que a gamificação pode aumentar o interesse dos alunos e melhorar a retenção do conteúdo, conforme discutido em diversas publicações da área.

\begin{figure}[h]
    \centering
    % \includegraphics[width=0.8\columnwidth]{placeholder} % Descomente e adicione sua imagem
    \caption{Exemplo de um diagrama de fluxo do jogo. Uma imagem pode ser inserida aqui para ilustrar um conceito. A pasta padrão para imagens é 'figuras/'.}
    \label{fig:fluxo_jogo}
\end{figure}

% ============================================================
% sections/02_problema.tex — Definição do Problema
% ============================================================

\section{Definição do Problema}

O ensino tradicional da tabela periódica baseia-se predominantemente em métodos
expositivos e mnemônicos, que exigem do estudante a memorização de um volume
considerável de informações desconexas do seu cotidiano. Essa abordagem frequentemente
resulta em aprendizado superficial e de curta duração \cite{gee2003videogames}.

O problema central que este trabalho busca endereçar pode ser formulado da
seguinte forma: \textit{como tornar o aprendizado dos elementos químicos da
tabela periódica mais engajador, significativo e eficaz por meio de um jogo
digital?}

% TODO: Adicionar referências a estudos sobre dificuldades no ensino de química
% TODO: Incluir dados de reprovação ou baixo desempenho em química no ensino médio

\subsection{Contexto Educacional}

Pesquisas recentes apontam que metodologias ativas de ensino, que colocam o
estudante como protagonista do processo de aprendizagem, produzem resultados
significativamente superiores em termos de retenção e compreensão do conteúdo
\cite{mayer2019serious}.

% Exemplo de equação (cálculo de eficiência de aprendizado — hipotético)
A eficiência de aprendizado $\eta$ pode ser modelada de forma simplificada como:

\begin{equation}
    \eta = \frac{C_r}{C_t} \times 100\%
    \label{eq:eficiencia}
\end{equation}

\noindent onde $C_r$ representa o conteúdo retido após determinado período e
$C_t$ o conteúdo total ensinado. A Eq.~\eqref{eq:eficiencia} ilustra que
metodologias que aumentam $C_r$ sem reduzir $C_t$ são intrinsecamente mais
eficientes.

\subsection{Lacuna Identificada}

Embora existam soluções digitais para o ensino de química, a maioria consiste
em simuladores ou quizzes que não exploram plenamente as mecânicas de engajamento
consolidadas pela indústria de jogos. O \textit{Alchemon} propõe preencher essa
lacuna ao incorporar sistemas de progressão, coleção e batalha diretamente
atrelados às propriedades dos elementos químicos.

% TODO: Realizar revisão sistemática de aplicativos existentes na área
% TODO: Comparar funcionalidades com solução proposta em tabela

% ============================================================
% sections/03_objetivos.tex — Objetivos
% ============================================================

\section{Objetivos}

\subsection{Objetivo Geral}

Desenvolver o \textit{Alchemon}, um jogo sério para plataformas digitais que
utilize mecânicas de coleção e batalha de criaturas baseadas nos elementos
químicos, com o objetivo de apoiar o ensino da tabela periódica no ensino médio.

\subsection{Objetivos Específicos}

Os objetivos específicos deste trabalho são:

\begin{enumerate}[i.]
    \item Realizar levantamento bibliográfico sobre jogos sérios e sua aplicação
          no ensino de ciências \cite{prensky2001digital};
    \item Definir o escopo, mecânicas e conteúdo educacional do jogo com base
          em princípios pedagógicos consolidados;
    \item Projetar a arquitetura do sistema e o design de personagens baseados
          nos elementos químicos;
    \item Implementar o protótipo funcional do jogo;
    \item Avaliar o protótipo por meio de testes com usuários e análise de
          métricas de aprendizagem.
\end{enumerate}

% TODO: Refinar objetivos específicos com orientador
% TODO: Verificar alinhamento com a ementa do curso

\subsection{Delimitação do Escopo}

O presente trabalho, por tratar-se de TCC I, abrange as etapas de pesquisa
bibliográfica, definição de requisitos, design do jogo e desenvolvimento do
protótipo inicial. A avaliação com usuários reais e a versão final do produto
serão conduzidas no âmbito do TCC II.

% Exemplo de tabela de objetivos vs. entregas
\begin{table}[htbp]
    \centering
    \caption{Mapeamento de Objetivos e Entregas — TCC I}
    \label{tab:objetivos_entregas}
    \begin{tabular}{@{}p{0.45\columnwidth}p{0.45\columnwidth}@{}}
        \toprule
        \textbf{Objetivo} & \textbf{Entrega} \\
        \midrule
        Revisão bibliográfica & Referencial teórico \\
        Design do jogo        & Documento de Design (GDD) \\
        Prototipação          & Protótipo jogável \\
        \bottomrule
    \end{tabular}
\end{table}

A Tabela~\ref{tab:objetivos_entregas} apresenta o mapeamento entre os objetivos
específicos e os artefatos entregues ao final deste trabalho.

% Chapter 4: Justification
\section{Justificativa}

A relevância deste projeto reside em três pilares principais: educacional, tecnológico e social.

Do ponto de vista educacional, este trabalho contribui com uma nova ferramenta didática que aborda um desafio persistente no ensino de química. A Tabela Periódica é a base para a compreensão de inúmeros conceitos químicos, e facilitar seu aprendizado pode ter um impacto positivo significativo no desempenho dos alunos na disciplina. A abordagem de aprendizagem baseada em jogos está alinhada com teorias construtivistas modernas, que defendem que o conhecimento é construído ativamente pelo aprendiz \cite{journal_example}.

Tecnologicamente, o projeto explora a aplicação de tecnologias de desenvolvimento de jogos para a criação de soluções educacionais, uma área em plena expansão. O desenvolvimento de "Alchemon" servirá como um estudo de caso prático sobre o design e a implementação de jogos sérios.

Socialmente, a democratização do acesso a ferramentas educacionais de qualidade por meio de dispositivos móveis pode ajudar a reduzir disparidades no aprendizado. Um jogo como "Alchemon" pode ser utilizado tanto em sala de aula, como um complemento às atividades do professor, quanto em casa, como uma forma de estudo autônomo.

% Chapter 5: Theoretical Foundation
\section{Fundamentação Teórica}

Esta seção apresenta a base teórica que sustenta o projeto, abrangendo os principais conceitos relacionados a jogos sérios, gamificação e teorias de aprendizagem aplicadas ao design de jogos.

\subsection{Jogos Sérios e Gamificação}
Jogos sérios são jogos projetados para um propósito primário que não é o puro entretenimento. Conforme discutido por autores como Miyamoto \cite{conference_example}, o design de jogos pode ser aplicado a diversos contextos, incluindo treinamento, simulação e educação. A gamificação, por sua vez, refere-se ao uso de elementos e mecânicas de jogos em contextos de não-jogo para aumentar o engajamento e a motivação dos usuários.

\subsection{Teorias da Aprendizagem}
O design de "Alchemon" será influenciado por teorias da aprendizagem como o Construtivismo e o Conectivismo. O Construtivismo, popularizado por Piaget, postula que o aprendizado é um processo ativo de construção de novos conhecimentos com base em experiências anteriores. Um jogo que permite a exploração e a experimentação, como o proposto, alinha-se perfeitamente com esta visão.

O design de objetos e interfaces também deve considerar os princípios de usabilidade e design centrado no usuário, para garantir que a ferramenta seja intuitiva e fácil de usar, como Norman descreve em seu livro seminal \cite{book_example}. A clareza da interface é crucial para não criar barreiras ao aprendizado.

\begin{figure}[h]
    \centering
    % \includegraphics[width=0.9\columnwidth]{placeholder2}
    \caption{Exemplo de um ciclo de feedback em um jogo sério, onde o jogador age, recebe feedback e aprende.}
    \label{fig:feedback_loop}
\end{figure}

% ============================================================
% sections/06_metodologia.tex — Metodologia
% ============================================================

\section{Metodologia}

Este trabalho adota uma abordagem de pesquisa aplicada com método de
desenvolvimento iterativo, combinando revisão bibliográfica sistemática,
design centrado no usuário e prototipação ágil \cite{mayer2019serious}.

\subsection{Classificação da Pesquisa}

Quanto à natureza: pesquisa aplicada, pois visa gerar um produto com
aplicação prática imediata. Quanto à abordagem: qualitativa e quantitativa
(mista). Quanto aos objetivos: exploratória e descritiva. Quanto aos
procedimentos: pesquisa bibliográfica e estudo de desenvolvimento.

% TODO: Fundamentar classificação metodológica com referência (ex: Gil, 2017)

\subsection{Etapas do Desenvolvimento}

O desenvolvimento do \textit{Alchemon} está organizado nas seguintes etapas:

\begin{enumerate}[1.]
    \item \textbf{Pesquisa bibliográfica}: revisão de literatura sobre jogos
          sérios, GBL e ensino de química;
    \item \textbf{Análise de requisitos}: levantamento das necessidades
          educacionais e técnicas;
    \item \textbf{Design do jogo}: elaboração do \textit{Game Design Document} (GDD);
    \item \textbf{Prototipação}: implementação do protótipo de baixa e alta fidelidade;
    \item \textbf{Avaliação} (TCC II): testes com usuários e análise de resultados.
\end{enumerate}

\subsection{Ferramentas e Tecnologias}

% TODO: Confirmar motor de jogo com orientador e atualizar esta subseção

\begin{table}[htbp]
    \centering
    \caption{Ferramentas Utilizadas no Desenvolvimento}
    \label{tab:ferramentas}
    \begin{tabular}{@{}lll@{}}
        \toprule
        \textbf{Ferramenta} & \textbf{Categoria} & \textbf{Versão} \\
        \midrule
        Unity               & Motor de Jogo      & 2022 LTS         \\
        Visual Studio Code  & Editor de Código   & 1.x              \\
        Blender             & Modelagem 3D       & 4.x              \\
        Git / GitHub        & Controle de Versão & ---              \\
        \bottomrule
    \end{tabular}
\end{table}

A Tabela~\ref{tab:ferramentas} lista as principais ferramentas adotadas no
projeto, selecionadas por serem gratuitas, amplamente utilizadas na indústria
e compatíveis com o ambiente de desenvolvimento do autor.

\subsection{Métricas de Avaliação}

A eficácia educacional do \textit{Alchemon} será mensurada por meio do
coeficiente de ganho normalizado de Hake:

\begin{equation}
    \langle g \rangle = \frac{\langle S_{pós} \rangle - \langle S_{pré} \rangle}
                             {100\% - \langle S_{pré} \rangle}
    \label{eq:hake}
\end{equation}

\noindent onde $\langle S_{pré} \rangle$ e $\langle S_{pós} \rangle$ representam
as médias dos escores percentuais em pré-teste e pós-teste, respectivamente.
Valores de $\langle g \rangle > 0{,}3$ são considerados ganhos satisfatórios
% NOTA: O uso de {,} em modo matemático requer \usepackage{icomma} ou siunitx
\cite{gee2003videogames}.

% TODO: Detalhar protocolo de coleta de dados e comitê de ética (se aplicável)

% Chapter 7: Game Planning
\section{Planejamento do Jogo "Alchemon"}

Esta seção detalha o conceito e o design do jogo "Alchemon".

\subsection{Conceito}
"Alchemon" é um RPG de batalha e coleta, onde os "monstros" são personificações dos elementos químicos ("Alchemons"). O jogador assume o papel de um jovem alquimista que explora um mundo de fantasia, captura Alchemons e os utiliza em batalhas. Cada Alchemon possui habilidades e atributos baseados nas propriedades reais do elemento químico que representa. Por exemplo, um Alchemon do tipo "Hélio" (gás nobre) pode ter habilidades defensivas que o tornam "inerte" a certos ataques.

\subsection{Mecânicas Principais}
\begin{itemize}
    \item 	extbf{Exploração:} O jogador navega por um mapa-múndi dividido em regiões que correspondem a grupos da Tabela Periódica (metais alcalinos, halogênios, etc.).
    \item 	extbf{Captura:} Ao encontrar um Alchemon selvagem, o jogador pode tentar capturá-lo respondendo a uma pergunta sobre o elemento correspondente (e.g., "Qual o número atômico do Carbono?").
    \item 	extbf{Batalha:} As batalhas ocorrem em turnos. A eficácia dos ataques depende das interações entre os tipos de Alchemons, que simulam reações químicas. Por exemplo, um ataque de um Alchemon do tipo "Sódio" pode ser super eficaz contra um do tipo "Cloro", formando "Sal".
\end{itemize}

O progresso no jogo está diretamente ligado ao aprendizado, incentivando o jogador a estudar a Tabela Periódica para se tornar mais forte no jogo, uma ideia inspirada por \cite{journal_example}.
A complexidade de uma batalha pode ser modelada pela seguinte relação:
\begin{equation}
    C = \sum_{i=1}^{n} P_i 	imes T_i
\end{equation}
onde \(C\) é a complexidade, \(P_i\) é a propriedade do elemento \(i\), e \(T_i\) é o tipo de interação.

% Chapter 8: Conclusion
\section{Conclusão e Trabalhos Futuros}

Este projeto de TCC I delineou a concepção e o planejamento do jogo sério "Alchemon", uma ferramenta destinada a apoiar o ensino da Tabela Periódica. Foi realizada a fundamentação teórica, a definição do problema, dos objetivos e da metodologia de desenvolvimento. O design do jogo foi estruturado para integrar de forma coesa os conteúdos pedagógicos com mecânicas de jogo engajadoras, seguindo as melhores práticas da área de aprendizagem baseada em jogos \cite{book_example, journal_example}.

O trabalho realizado até aqui corresponde às fases iniciais do ciclo de desenvolvimento, resultando em um plano sólido para a implementação do protótipo. Acreditamos que a abordagem proposta tem grande potencial para criar uma experiência de aprendizagem positiva e eficaz.

\subsection{Trabalhos Futuros}
A continuação deste projeto, a ser realizada no TCC II, envolverá as seguintes etapas:
\begin{itemize}
    \item Conclusão do desenvolvimento do protótipo funcional do "Alchemon".
    \item Realização de testes de usabilidade para refinar a interface e a experiência do usuário.
    \item Implementação de um estudo de caso em um ambiente educacional real, com a participação de estudantes do ensino médio.
    \item Coleta e análise de dados quantitativos e qualitativos para avaliar o impacto do jogo na motivação e no desempenho acadêmico dos alunos.
    \item Documentação dos resultados e redação das conclusões finais da pesquisa.
\end{itemize}

Espera-se que, ao final do projeto completo, o "Alchemon" não seja apenas um produto final, mas também uma contribuição validada para o campo da tecnologia educacional, como sugerido em conferências sobre inovação \cite{conference_example}.


% ------------------------------------------------------------
% REFERÊNCIAS BIBLIOGRÁFICAS
% OBRIGATÓRIO: bibliographystyle + bibliography para bibtex
% ------------------------------------------------------------
\bibliographystyle{IEEEtran}
\bibliography{referencias}

\end{document}
